\section{Supporting mass-readout arguments}
For completeness, the original mathematical arguments for the conditional mass-readout statements are reproduced below. Their hypotheses refer to the comparison domains stated in the text; they do not identify a native physical readout.
Let $x_i=\sqrt{m_i}\geq0$, $x\ne0$, and write
\[
 A=\frac{x_0+x_1+x_2}{3}>0,\qquad
 u=A(1,1,1),\quad d=x-u,\quad \sum_i d_i=0.
\]
The cyclic permutation preserves the one-dimensional symmetric line
and its two-dimensional orthogonal plane. It preserves their norms
separately; it does not prescribe their ratio.
\begin{proposition}[Koide as an equal norm condition]
\label{v480:koide-equal-norm}
For these nonnegative principal roots,
\[
 K=\frac{\sum_i x_i^2}{(\sum_i x_i)^2}
   =\frac13+\frac{\|d\|^2}{9A^2}.
\]
Consequently $K=2/3$ if and only if
$\|d\|^2=\|u\|^2=3A^2$.
\end{proposition}
\begin{proof}
Since $u\perp d$, $\|x\|^2=3A^2+\|d\|^2$, while
$(\sum_i x_i)^2=9A^2$. Division proves both assertions.
\end{proof}
\begin{proposition}[A sufficient HC shaped mass root realization]
\label{v480:koide-hc-mixing}
In a declared real three-root representation, let
$e=(1,1,1)/\sqrt3$, let $J_m$ be an orthogonal involution, and put
$C_m=(I+J_m)/2$. Suppose $x=\alpha C_m e$ is nonzero and
componentwise nonnegative, with $\alpha>0$. For
$h=\langle e,J_m e\rangle$, one has $h>-1$ and
\[
 K=\frac{2}{3(1+h)}.
\]
Hence Koide holds precisely when $\langle e,J_m e\rangle=0$.
\end{proposition}
\begin{proof}
An orthogonal involution is self-adjoint, so $C_m$ is an
orthogonal projector. Thus
$\|C_m e\|^2=\langle e,C_m e\rangle=(1+h)/2>0$ and
$(\sum_i x_i)^2=3\alpha^2(1+h)^2/4$. Dividing
$\|x\|^2=\alpha^2(1+h)/2$ by this denominator gives the formula.
\end{proof}
\begin{proposition}[Radius and phase in the three root parametrization]
\label{v480:koide-radius-phase}
Every nonzero nonnegative three-root vector admits
\[
 x_k=A\bigl(1+r\cos(\theta+2\pi k/3)\bigr),\qquad r\geq0,
\]
with $K=1/3+r^2/6$. The condition $K=2/3$ is precisely $r=\sqrt2$;
it does not fix $A$ or $\theta$.
\end{proposition}
\begin{proof}
The vectors $(\cos(2\pi k/3))_{k=0}^2$ and
$(\sin(2\pi k/3))_{k=0}^2$ form an orthogonal basis of the
zero-sum plane, each with squared norm $3/2$. Expressing $d/A$
in that basis gives coefficients $r\cos\theta$ and
$-r\sin\theta$. Therefore $\|d\|^2=(3/2)A^2r^2$ and the
previous proposition gives the formula for $K$. If $r=0$ the
phase is arbitrary.
\end{proof}
\begin{proposition}[Cyclic equivariant linear maps preserve symmetric inputs]
\label{v480:koide-equivariant-obstruction}
Let $P$ cyclically permute three real coordinates. If a linear map
$T$ satisfies $TP=PT$, then $T(A(1,1,1))$ remains on the symmetric
line. It cannot produce a nondegenerate root hierarchy from that input.
\end{proposition}
\begin{proof}
$Pu=u$ implies $P(Tu)=T(Pu)=Tu$. The fixed space of $P$ is
exactly the line spanned by $(1,1,1)$.
\end{proof}
\begin{proposition}[descent]
\label{v480:hc-mass-descent}
Let $H$ be a real Hilbert space, $J$ an orthogonal involution, and
$B:H\to\mathbb R^3$ a bounded coisometry, $BB^*=I_3$. If
$\ker B$ is $J$-invariant, then
\[
 J_m=BJB^*,\qquad J_m^*=J_m,\qquad J_m^2=I_3,
 \qquad BC=C_mB,\quad C_m=(I_3+J_m)/2.
\]
Conversely, any relation $BJ=J_mB$ requires $\ker B$ to be
$J$-invariant. The descended operator is unique because $B$ is onto.
\end{proposition}
\begin{proof}
Put $E=B^*B$, the orthogonal projection onto $(\ker B)^\perp$.
Since $J$ is self-adjoint and preserves $\ker B$, it also preserves
its orthogonal complement, so $JE=EJ$. Thus $J_m^*=J_m$ and
$J_m^2=BJEJB^*=BEB^*=I_3$. Also
$J_mB=BJE=BJ$, since $J$ preserves $\ker B$ and $B$ kills it.
The identity for $C$ follows by averaging. Conversely,
$Bk=0$ implies $BJk=J_mBk=0$. Surjectivity gives uniqueness.
\end{proof}
\begin{proposition}[odd obstruction]
\label{v480:hc-mass-odd-obstruction}
If $Jf=-f$, then $Cf=0$. Consequently any homogeneous linear
root readout $x=BCf$ is zero. Any readout $F(Cf)$ with $F(0)=0$
has the same obstruction on this input lane. If instead a seed is
$J$-fixed and the descended linear route has $Bf_0=e$, then
$J_me=e$ and its nonzero root output has $K=1/3$.
\end{proposition}
\begin{proof}
The first statement follows from $Cf=(f+Jf)/2=0$.
For the second, $J_me=BJf_0=Bf_0=e$, so $C_me=e$ and
all three roots are equal.
\end{proof}
\begin{proposition}[equivariant obstruction]
\label{v480:hc-mass-equivariant-obstruction}
Let $U$ preserve an admitted state set and let $P$ cyclically permute
the three root coordinates. For any readout $F$ satisfying
$F(Uf)=PF(f)$, an input $Uf_0=f_0$ has
$F(f_0)=a(1,1,1)$. If this is a nonzero nonnegative root vector,
then $K=1/3$. No linearity assumption on $F$ is needed.
\end{proposition}
\begin{proof}
Equivariance gives $PF(f_0)=F(Uf_0)=F(f_0)$. The fixed vectors
of the three-cycle are exactly the constant triples.
\end{proof}
\begin{proposition}[triole bound]
\label{v480:hc-mass-triole-bound}
Let $z_0+z_1+z_2=0$ in a normed vector space, not all $z_i=0$.
Under the specific identification $m_i=\kappa\|z_i\|^2$ with a
common $\kappa>0$, the principal-root ratio satisfies
\[
 \frac13\leq K\leq\frac12.
\]
In particular it cannot equal $2/3$.
\end{proposition}
\begin{proof}
Put $r_i=\|z_i\|$ and $s=\sum_i r_i>0$. The triangle inequality
gives $r_i\leq s-r_i$, hence $r_i\leq s/2$. Therefore
$\sum_i r_i^2\leq(s/2)\sum_i r_i=s^2/2$.
Cauchy--Schwarz gives the lower bound, and the common mass scale
cancels from $K$.
\end{proof}
\begin{proposition}[fit freedom]
\label{v480:hc-mass-fit-freedom}
For every nonzero nonnegative triple $x$, there exist an orthogonal
involution $J_x$ and $\alpha_x>0$ such that
$x=\alpha_x(I+J_x)e/2$. Thus existence of an HC-shaped operator
chosen after the masses is not a prediction of Koide or a hierarchy.
\end{proposition}
\begin{proof}
Set $C_x=xx^{\mathsf T}/\|x\|^2$ and $J_x=2C_x-I_3$.
This is an orthogonal involution. Since $s=\sum_i x_i>0$,
$C_xe=sx/(\sqrt3\|x\|^2)$, so choose
$\alpha_x=\sqrt3\|x\|^2/s$.
\end{proof}
